% Do not modify these
\documentclass[fleqn,10pt]{wlscirep}
\usepackage[utf8]{inputenc}
\usepackage[T1]{fontenc}



% -- Insert any custom LaTeX packages here --

% \package{natbib} % <-- Required for the Chicago citation style
% \package{apacite} % <-- Required for the APA citation style
% If you decide to use one of the styles above, remember to change the \bibliographystyle{} at the bottom of the document too!

\usepackage{listings} % <-- Required if you want to display program source code in your paper.

\usepackage{hyperref}
\hypersetup{
    colorlinks=true,
    linkcolor=blue,
    filecolor=magenta,      
    urlcolor=cyan,
}
 
\urlstyle{same}

% -- End of custom LaTeX packages --


% Fill in your title
\title{Exploring Bitcoin }

% Do not modify the author tag below, just let it be blank
\author{}

% Fill in assignment abstract
\begin{abstract}
In this assignment we parsed a dataset, in the form of XMLs, through a Hadoop- and Pig apache. We learned how to write MapReduce jobs and Pig scripts, and how effective it is on our dataset based on XML files from Bitcoin stack exchange. 

\end{abstract}


% Do not modify the following two lines
\begin{document}
% define the infrastructure
\ExplSyntaxOn
\keys_define:nn { coverpage }
 {
  assignment_title .tl_set:N 	= \hk_variable_assignment_title_tl,
  student_name  .tl_set:N 		= \hk_variable_student_name_tl,
  student_number     .tl_set:N 	= \hk_variable_student_number_tl,
  course_name .tl_set:N 			= \hk_variable_course_name_tl,
  course_code .tl_set:N			= \hk_variable_course_code_tl,
  due_date .tl_set:N 			= \hk_variable_due_date_tl,
  master_of .tl_set:N 			= \hk_variable_master_of_tl,
  group_size .tl_set:N 			= \hk_variable_group_size_tl,
 }
\cs_new_protected:Nn \use_coverpage:n
 {
  \keys_set:nn { coverpage } { #1 }
  \coverpage
 }
\cs_set_eq:NN \makecoverpage \use_coverpage:n
\cs_new:Nn \use_variable:n
 {
  \tl_use:c { hk_variable_#1_tl }
 }
\cs_set_eq:NN \use \use_variable:n
\ExplSyntaxOff



% Command for making the cover page

\newcommand{\coverpage}{%
\begin{titlepage}
\begin{center}
\begin{figure}[ht]
\centering
\includegraphics[width=0.5\linewidth]{images/hk-logo.pdf}
\end{figure}
\huge{Course Assignment}\\
\LARGE{Master of \use{master_of}}\\
\Large{Department of Technology}\\

\vskip10pt%
\begin{table}[ht]
\centering
\Large
\setlength{\tabcolsep}{0.5em}
\setlength\extrarowheight{1.2em}
\begin{tabularx}{0.8\linewidth}{lX}
\toprule
 Assignment title 	& \use{assignment_title} \\ \hline
 Course code 		& \use{course_code} \\ \hline
 Course name 		& \use{course_name} \\ \hline
 Due date			& \use{due_date} \\
 \bottomrule
\end{tabularx}%
\end{table}

\vskip10pt
\textbf{\Large{Declaration:}}\\

\if\use{group_size}1
Through the submission of this assignment, I hereby declare\\ that this report is the result of my own work, and that all\\ sources have been properly cited to throughout the text.
\vskip10pt%
\begin{table}[ht]
\centering
\Large
\setlength{\tabcolsep}{0.5em}
\setlength\extrarowheight{1em}
\begin{tabularx}{0.8\linewidth}{lX}
\toprule
 Student name 		& \use{student_name} \\ \hline
 Student ID number 	& \use{student_number} \\ 
 \bottomrule
\end{tabularx}%
\end{table}

\else 

Through the submission of this assignment, we hereby declare\\ that this report is the result of our own work, and that all\\ sources have been properly cited to throughout the text.
\begin{table}[ht]
\centering
\Large
\setlength{\tabcolsep}{0.5em}
\setlength\extrarowheight{1em}
\begin{tabularx}{0.8\linewidth}{lX}
\toprule
 Names of students 		& \use{student_name} \\ \hline
 Student ID numbers 	& \use{student_number} \\ 
 \bottomrule
\end{tabularx}%
\end{table}
\fi




\end{center}
\end{titlepage}
}


% Insert data for the hand-in's cover page
\makecoverpage{
	master_of 		 = \par{Applied Computer Science},  % Use either: Applied Computer Science | Human-Computer Interaction
	assignment_title = \par{ Exploring Bitcoin} ,  % Title of your assignment
	course_code    	 = \par{MA120},  % Course code (ex. MA110)
	course_name      = \par{Big Data},  % Course name (ex. Systems Development)
	due_date		 = \par{18 Oktober},  % Due date
	student_name     = \par{Theodore Midtbø Alstad ; Howie Chen},  % Your name (or names, if group – separate names with ; semicolon)
	student_number   = \par{865317 ; 866354},  % Your student ID number (or numbers, if group – separate ID numbers with ; semicolon)
	group_size		 = 2, % Number of group members (used for the declaration text)
}


% Do not modify the following two lines
\flushbottom
\maketitle


% --INTRODUCTION--
\section*{Introduction}
We choose to work together because both of us have python background, therefor we choose python as main programming language. We explored Bitcoin as given dataset from \url{archive.org/download/stackexchange/bitcoin.stackexchange.com.7z}  through Apaches: Hadoop and Pig. 

% --FUNCTIOS %-- 
\section*{Main functions}

% --TASK1--
\section*{Task 1 Warmup}
%This part of the task is about get to know how hadoop- and pig apache works and how to parse XML  files throught python code. 
\subsection*{1a) WordCount}
\textbf{Assumption}: Count the words in body of questions PostTypeID="1"  in the \textit{Posts.xml}  file. The result should be how many times a word occur in the body of questions.\\ \\
\textbf{Implementation}Hello hello \\ \\
\textbf{Notes/Reflection} bye bye  \\ \\

\subsection*{1b) Unique words}
\textbf{Assumption}: Write a MapReduce job where the result should be unique words in the titles PostTypeID="1" in the \textit{Posts.xml} File. \\ \\
\textbf{Implementation}  \\ \\
\textbf{Notes/Reflection}

\subsection*{1c) MoreThan10}
\textbf{Assumption}: Write a simple python code to check how many words there is in the title in the \textit{Posts.xml}. The result should output how many titles are longer than 10 words.  \\ \\
\textbf{Implementation}  \\ \\
\textbf{Notes/Reflection}

\subsection*{1d) Stopwords}
\textbf{Assumption}: Write a simple python code based on task 1a to exclude \href{https://raw.githubusercontent.com/naimdjon/stopwords/master/stopwords.txt}{stopwords} from body of questions PostTypeID="1" in the \textit{Posts.xml}. The output should be textfile without any stopwords. \\ \\
\textbf{Implementation}  \\ \\
\textbf{Notes/Reflection}

\subsection*{1e) Pig top 10}
\textbf{Assumption}: Write a pig script to select top 10 listed words after removing the stopwords from \textit{Posts.xml}. The output should print out top 10 listed words.\\ \\
\textbf{Implementation}  \\ \\
\textbf{Notes/Reflection}

\subsection*{1f) Tags}
\textbf{Assumption}: Write a MapReduce job to create a dictionary over unique tags in  \textit{Posts.xml}. \\ \\
\textbf{Implementation}  \\ \\
\textbf{Notes/Reflection}

% --TASK2--
\section*{Task 2 Discover}
%This part of the task is about to looking through several 

\subsection*{2a Counting) }
\textbf{Assumption}: Write a simple python code to list how many unique users in \textit{Users.xml}. \\ \\
\textbf{Implementation}  \\ \\
\textbf{Notes/Reflection}

\subsection*{2b Unique users) }
\textbf{Assumption}:  \\ \\
\textbf{Implementation}  \\ \\
\textbf{Notes/Reflection}

\subsection*{2c Top miners) }
\textbf{Assumption}:  \\ \\
\textbf{Implementation}  \\ \\
\textbf{Notes/Reflection}

\subsection*{2d Top questions) }
\textbf{Assumption}:  \\ \\
\textbf{Implementation}  \\ \\
\textbf{Notes/Reflection}

\subsection*{2e Favorite questions) }
\textbf{Assumption}:  \\ \\
\textbf{Implementation}  \\ \\
\textbf{Notes/Reflection}

\subsection*{2f Average answers) }
\textbf{Assumption}:  \\ \\
\textbf{Implementation}  \\ \\
\textbf{Notes/Reflection}

\subsection*{2g Countries) }
\textbf{Assumption}:  \\ \\
\textbf{Implementation}  \\ \\
\textbf{Notes/Reflection}

\subsection*{2h Names) }
\textbf{Assumption}:  \\ \\
\textbf{Implementation}  \\ \\
\textbf{Notes/Reflection}

\subsection*{2i) Answers }
\textbf{Assumption}:  \\ \\
\textbf{Implementation}  \\ \\
\textbf{Notes/Reflection}

% --TASK3--
\section*{Task 3 Numbers}

\subsection*{3a) Bigram }
\textbf{Assumption}:  \\ \\
\textbf{Implementation}  \\ \\
\textbf{Notes/Reflection}

\subsection*{3b) Trigram }
\textbf{Assumption}:  \\ \\
\textbf{Implementation}  \\ \\
\textbf{Notes/Reflection}

\subsection*{3c) Combiner }
\textbf{Assumption}:  \\ \\
\textbf{Implementation}  \\ \\
\textbf{Notes/Reflection}

\subsection*{3d) Useless}
\textbf{Assumption}:  \\ \\
\textbf{Implementation}  \\ \\
\textbf{Notes/Reflection}

% --TASK4--
\section*{Task 4 Search engine}
\subsection*{4a) Title index}
\textbf{Assumption}:  \\ \\
\textbf{Implementation}  \\ \\
\textbf{Notes/Reflection}

\section*{Conclusion}
% Do not modify this last lines
\end{document}